\section{What a Non-Null Field Promises}
\label{sec:ch04-what-a-non-null-field-promises}

\index{nullability!comes from the C\# annotation|(}%
Nobody has typed an exclamation mark into this book's schema. Every one on a
field you wrote came from a C\# annotation: \code{Nullable} has been enabled in
the project file since chapter~\ref{ch:hot-chocolate-zero}, and Hot Chocolate
reads it straight across.

\begin{minted}{text}
string Title           ->  title: String!
string? Abstract       ->  abstract: String
Task<Speaker?>         ->  speaker: Speaker
Task<Speaker>          ->  speaker: Speaker!
\end{minted}

\index{nullability!the question mark is API design}%
So a question mark in a C\# signature is an API decision made in a file that
does not look like one. Chapter~\ref{ch:data-without-n-plus-1} chose
\code{Task<Speaker?>} because a DataLoader answers a missing key with null and
the type had to allow it. That was a local choice about a return value, and it
was also a permanent promise to every client the graph will ever have.

\index{non-null!what it costs to be wrong}%
Permanent because there is no safe direction out of it. Adding an exclamation
mark later breaks every client that handled null, whose null branch is now
dead code guarding against a case the schema says cannot occur. Removing one
breaks every client that trusted the field and never wrote that branch at all.
The nullability of a field is a decision you make once.

\subsection{One Missing Row}

\index{non-null!and a missing row}%
So take the promise chapter~\ref{ch:data-without-n-plus-1} did not make, and
make it. One character comes off the speaker resolver, and session~4 in
\code{ConferenceData} is pointed at speaker~99, who does not exist.

\begin{minted}[highlightlines={12}]{csharp}
using GreenDonut.Data;
using Microsoft.EntityFrameworkCore;

namespace Sessions;

[ObjectType<Session>]
public static partial class SessionType
{
    /// <summary>
    /// The person giving this session.
    /// </summary>
    public static async Task<Speaker> GetSpeakerAsync(
        [Parent(nameof(Session.SpeakerId))] Session session,
        ISpeakerByIdDataLoader speakerById,
        CancellationToken ct)
        => await speakerById.LoadAsync(session.SpeakerId, ct);

    /// <summary>
    /// Loads one session by the identifier inside a global node id.
    /// </summary>
    [NodeResolver]
    public static async Task<Session?> GetSessionAsync(
        int id,
        ConferenceContext db,
        QueryContext<Session> query,
        CancellationToken ct)
        => await db.Sessions
            .Where(s => s.Id == id)
            .With(query)
            .FirstOrDefaultAsync(ct);
}
\end{minted}

\index{CS8603!the compiler saw it first}%
The build still succeeds, and it does not go quietly:

\begin{minted}{text}
warning CS8603: Possible null reference return.
\end{minted}

The method promises a \code{Speaker}, the loader can hand it back a null, and
the compiler can see the gap before a single request has been sent. The schema
now carries \code{speaker: Speaker!}, so the same gap is in the contract.

Ask for the page.

\begin{minted}{text}
{"query":"{ sessions(first: 10) { nodes { title speaker { name } } } }"}
\end{minted}

\begin{minted}[fontsize=\footnotesize]{json}
{"errors":[{"message":"Cannot return null for non-nullable field.",
"path":["sessions","nodes",3,"speaker"],"extensions":{"code":"HC0018"}}],
"data":{"sessions":{"nodes":null}}}
\end{minted}

\index{blast radius!four sessions for one bad row}%
Four sessions asked for, none returned. Sessions one, two and three resolved
perfectly well and are gone with the fourth.

\index{status 200!with an empty list}%
Status 200, incidentally, and one error rather than one per lost session: the
count follows the position that failed, not the rows that vanished.

\subsection{Why It Climbs}

\index{null propagation!the specification's rule}%
The specification describes this as a walk upward rather than as a
rule about lists~\autocite{graphqlspec2025}:

\begin{quote}
Since Non-Null response positions cannot be null, execution errors are
propagated to be handled by the parent response position. If the parent
response position may be null then it resolves to null, otherwise if it is a
Non-Null type, the execution error is further propagated to its parent
response position.
\end{quote}

\index{null propagation!through a list of non-null items}%
Walk it on this response. \code{speaker} is \code{Speaker!} and cannot hold the
null, so the error goes to its parent, which is the \code{Session} at index
three. \code{nodes} is typed \code{[Session!]}, so the items in it cannot be
null either, and the error keeps going. The parent of the item is the list, and
the list itself is nullable, so it stops there and \code{nodes} resolves to
null. Figure~\ref{fig:ch04-null-climbs} is that walk, drawn over a selection
that asks for \code{pageInfo} as well, because where the climb does not reach
is half of what the picture is for.

\begin{figure}
  \centering
  % One response, drawn as its tree of positions: a null raised at the speaker
% of session index 3 climbs until it reaches a position whose type is allowed
% to hold it.
%
% Same stroke weight, arrowhead and font as figures/tikz/ch01-one-field.tex,
% the file that fixes this book's idiom: 0.4pt strokes, the Straight Barb
% arrowhead, sans-serif scriptsize labels, square corners, no fills, black on
% white. That figure and ch03-one-batch.tex are timelines because their
% subject is order. This one is a shape, so it is a tree; the idiom is the
% line weight and the labelling, not the axis.
%
% Solid box: a position that is present in the response. That includes nodes,
% whose key survives even though its value is null. Dashed box: a position
% that does not exist in the response at all, because the list that would
% have held it is gone. Each box on the path carries its own GraphQL type,
% because the nullability of that type is the only thing deciding whether the
% climb stops there.
%
% The two solid arrows are the climb, drawn along the tree edges they run up.
% The query's title and name fields are left out: neither changes where the
% climb stops.
%
% Bare tikzpicture (house style): the figure environment, caption and label
% live at the call site in the section file.
\begin{tikzpicture}[
  x=1mm, y=1mm,
  every node/.append style={font=\sffamily\scriptsize},
  posbox/.style={draw, line width=0.4pt, minimum width=20mm,
    minimum height=9mm, align=center, inner sep=1pt},
  lostbox/.style={posbox, dashed},
  link/.style={draw, line width=0.4pt},
  lostlink/.style={draw, line width=0.4pt, dashed},
  climb/.style={draw, line width=0.4pt,
    -{Straight Barb[length=1.4mm, width=1.6mm]}},
  span/.style={draw, line width=0.4pt},
  spanlabel/.style={anchor=north, align=center, inner sep=2pt},
  note/.style={anchor=west, align=left, inner sep=2pt},
  haltnote/.style={anchor=east, align=right, text width=30mm, inner sep=2pt},
]

% --------------------------------------------------------------- positions
\node[posbox] (root)  at (83,78) {sessions};

\node[posbox] (nodes) at (46,54) {nodes\\{}[Session!]};
\node[posbox] (page)  at (120,54) {pageInfo};

\node[lostbox] (i0) at (10,30) {nodes[0]};
\node[lostbox] (i1) at (32,30) {nodes[1]};
\node[lostbox] (i2) at (54,30) {nodes[2]};
\node[lostbox] (i3) at (82,30) {nodes[3]\\Session!};

\node[lostbox] (speaker) at (82,6)  {speaker\\Speaker!};
\node[posbox]  (hasnext) at (120,30) {hasNextPage};

% ---------------------------------------------------------- tree structure
\draw[link] (root.south) -- (83,66) -- (46,66) -- (nodes.north);
\draw[link] (root.south) -- (83,66) -- (120,66) -- (page.north);
\draw[link] (page.south) -- (hasnext.north);

\draw[lostlink] (nodes.south) -- (46,42) -- (10,42) -- (i0.north);
\draw[lostlink] (nodes.south) -- (46,42) -- (32,42) -- (i1.north);
\draw[lostlink] (nodes.south) -- (46,42) -- (54,42) -- (i2.north);

% ------------------------------------------------------------- the climb
\draw[climb] (speaker.north) -- (i3.south);
\draw[climb] (i3.north) -- (82,42) -- (46,42) -- (nodes.south);

% ------------------------------------------------------------------ notes
% Each note sits on a row of its own. The types are carried inside the boxes
% rather than in side labels, because a label beside nodes[3] runs into the
% pageInfo branch at this width.
\node[note] at (94,6)
  {the resolver returned null here,\\and Speaker! cannot hold it};
\node[haltnote] at (32,54)
  {the first position that may be null, so the climb stops and nodes is
   null};
\node[spanlabel] at (120,24) {untouched};

% ------------------------------------------------------- collateral loss
\draw[span] (0,23) -- (0,19) -- (64,19) -- (64,23);
\node[spanlabel] at (32,17)
  {these three resolved correctly and are gone anyway};

% ----------------------------------------------------------------- legend
\node[posbox, minimum width=7mm, minimum height=4mm] at (4,-2) {};
\node[note] at (9,-2) {present in the response};

\node[lostbox, minimum width=7mm, minimum height=4mm] at (62,-2) {};
\node[note] at (67,-2) {not in the response at all};

\end{tikzpicture}

  \caption{The same response as a tree of positions, with the three positions
  on the path carrying their own types. The failure starts at the speaker of
  session four and climbs until it reaches one that is allowed to be null. Two
  on the way up are not, which is why three sessions that resolved correctly
  are lost with the fourth, and why the damage stops at \code{nodes} rather
  than reaching \code{sessions}. \code{pageInfo} is on the other branch and
  never learns anything happened.}
  \label{fig:ch04-null-climbs}
\end{figure}

\index{null propagation!stops at the first nullable position}%
The list is the position that absorbs it, and \code{[Session!]} is worth
stopping on because it reads like a stronger guarantee than \code{[Session]}
and is the shape I write by reflex. It is stronger, and the strength is what
does the damage: a list that may not contain a null has nowhere to put one, so
it goes null entirely rather than shortening by an element.

\index{blast radius!bounded by the nearest nullable position}%
Three more requests bound the damage. Ask for \code{pageInfo} alongside and it
comes back intact, because it is a sibling of \code{nodes} under
\code{sessions} and the climb never reaches their common parent. Ask through
\code{edges} instead of \code{nodes} and the same thing happens one level
deeper, with a path of \code{sessions}, \code{edges}, 3, \code{node},
\code{speaker}: the edge's \code{node} is non-null too, so \code{edges} is the
position that absorbs it. And ask for the orphaned session on its own through
\code{sessionById}, whose type is nullable, and the whole field is null with a
path two entries long. The blast radius is the distance from where the failure
happened to the nearest position that may be null, and nothing about the
failure itself sets it.

\index{selection set!decides whether the failure is visible}%
A fourth request is the uncomfortable one. Ask the same broken service for
\code{sessions \{ nodes \{ title \} \}}, without the speaker, and all four
titles come back with no error key anywhere. The schema is exactly as wrong.
Whether a given client ever sees it depends on what that client selected, so
this is a defect that arrives in production on the day somebody adds a field to
a query.

\subsection{The Rule I Use}

\index{nullability!the rule}%
Mark a field non-null when you can guarantee it from data you own, in the same
process, with a constraint behind it. \code{Session.title} is non-null because
the column is not nullable and the row cannot exist without one. Everything
else is nullable.

\index{nullability!a field that leaves the process}%
The load-bearing half is \enquote{data you own, in the same process}, and this
section has just measured why. \code{speaker} reads through a DataLoader whose
answer depends on a row in another table matching a foreign key nothing
enforces, and that is already too far. By
chapter~\ref{ch:second-third-subgraph} the same field is answered by a
different service over the network, where a non-null field is a promise about
another team's uptime written into your schema.

\index{non-null!where it is still right}%
None of which argues for a schema with no exclamation marks. Make everything
nullable and every field of every client grows a null check, paid on every
line of client code forever to insure against failures that mostly cannot
happen. Non-null on \code{id}, on a required column, on the fields of a type
you constructed yourself: keep all of it. What decides the field is who owns
the value, and how cautious you feel is not a factor.

\index{nullability!where the rest of this belongs}%
Chapter~\ref{ch:null-blast-radius} is where it gets expensive, because a
subgraph that fails is a whole branch of the response failing at once and the
walk in figure~\ref{fig:ch04-null-climbs} can reach \code{data} itself. It
also covers what the specification working group is doing about the fact that
the only tool here is a type modifier.
\index{nullability!comes from the C\# annotation|)}%
